% Generated draft: requires mathematical and source review.
\documentclass[11pt]{article}
\usepackage[margin=25mm]{geometry}
\usepackage{fontspec}
\setmainfont{Latin Modern Roman}
\usepackage{amsmath,amssymb,mathtools,unicode-math}
\setmathfont{Latin Modern Math}
\usepackage{enumitem,needspace,xurl,hyperref}
\setlength{\parindent}{0pt}
\setlength{\parskip}{5pt}
\setlength{\emergencystretch}{3em}
\newcommand{\N}{\mathbb{N}}
\newcommand{\Z}{\mathbb{Z}}
\newcommand{\Q}{\mathbb{Q}}
\newcommand{\R}{\mathbb{R}}
\newcommand{\C}{\mathbb{C}}
\newcommand{\F}{\mathbb{F}}
\newcommand{\A}{\mathbb{A}}
\newcommand{\PP}{\mathbb{P}}
\newcommand{\E}{\mathbb{E}}
\newcommand{\eps}{\varepsilon}
\newcommand{\dd}{\,\mathrm{d}}
\begin{document}
\section*{UnsolvedMath ID: 800317}
\textbf{Source problem code:} catalogue entry
\subsection{Definitions and mathematical statement}
Under the usual convention that graphs are finite, simple, and undirected, let \([n]=\{1,\ldots,n\}\) for each positive integer \(n\). A graph on \([n]\) has an edge set \(E\subseteq\binom{[n]}{2}\). Two such graphs are isomorphic if a permutation of \([n]\) carries one edge set to the other. Define
\[
 a_n=\left|\left\{E:E\subseteq\binom{[n]}{2}\right\}/\mathfrak S_n\right|,
\]
where \(\mathfrak S_n\) acts by permuting vertices. Thus \(a_n\) counts unlabeled graphs on \(n\) vertices under this convention.

The supplied conjecture, attributed to Wilf and listed as Conjecture 1 in Pak's 2018 source, asserts:
\begin{quote}
The sequence \(\{a_n\}\) has no Wilfian formula of type \(\mathrm{(W1)}\).
\end{quote}
The supplied record does not define a Wilfian formula of type \(\mathrm{(W1)}\), including its computational model or complexity requirement. Consequently, this assertion retains an undefined source term and is not yet a complete formal computational statement. The graph convention above is a conventional interpretation requiring confirmation against the source; the record also does not specify whether the sequence includes \(n=0\).
\subsection{Short English statement}
The conjecture attributed to Wilf says that the sequence counting unlabeled graphs has no Wilfian formula of type \(\mathrm{(W1)}\). The supplied excerpt omits the definition of that formula type, so a precise complexity statement requires further source information.
\subsection{Sources}
\begin{itemize}
\item Local UnsolvedMath registry, permanent ID \(800317\), source problem code \texttt{catalogue entry}, in the AMR Open Problem Lists collection. The supplied record labels the problem open; this status is not independently verified.
\item Pak, \emph{Complexity Problems in Enumerative Combinatorics} (2018), Conjecture 1. Supplied source URL: \url{https://arxiv.org/abs/1803.06636}. The record reports access on 2026-10-16 and extraction by \texttt{arxiv\_latex}; these are supplied provenance details, not a claim of independent source access.
\item Attribution/context: Wilf. The excerpt refers to an original paper labeled [Wilf], but supplies no title, publication details, or separate URL for that reference. The registry lists Pak as proposer and 2018 as the proposed year; these indexing fields do not resolve the original Wilf attribution or date.
\end{itemize}
\end{document}
